% =====================================================================
%  definitions.tex -- angepasste Kopie von Vorlagen/definitions.tex
%  (MCI-Vorlage V3.0, Drauschke/Teschl/Mehrle)
%
%  Layout, Farben, Überschriften, Abstände und Verzeichnisnamen sind
%  unverändert übernommen. Geändert wurde nur, was mit pdflatex /
%  aktuellem TeX Live nicht mehr kompiliert:
%    - "german"           -> babel/ngerman + UTF-8 (Umlaute direkt tippen)
%    - graphicx[dvips]    -> graphicx (pdflatex kann PNG/PDF einbinden)
%    - pst-all, wallpaper, pgfplots, units, subfigure entfernt
%      (Konflikte \unit/\ang, nur für latex+dvips gedacht, nicht benötigt)
%    - PSTricks-Makros (Netzwerk-/Blocksymbole) entfernt
%  Ergänzt: booktabs, listings, \todo für offene Stellen.
% =====================================================================

\ifthenelse{\boolean{english}}
  {\usepackage[english]{babel}}
  {\usepackage[ngerman]{babel}}
\usepackage[utf8]{inputenc}
\usepackage[T1]{fontenc}
\usepackage{lmodern}
\usepackage{graphicx}
\usepackage[tmargin=1in,bmargin=1in,lmargin=1.25in,rmargin=1.25in]{geometry}
\usepackage{titlesec}
\usepackage{xcolor}
\usepackage[overload]{textcase}
\definecolor{MSBlue}{rgb}{.204,.353,.541}
\definecolor{MSLightBlue}{rgb}{.31,.506,.741}
% Set formats for each heading level (Abschnitte klein geschrieben wie in der Vorlage)
\titleformat{\section}[hang]{\rmfamily\bfseries\huge\color{MSBlue}}{\thesection}{1em}{\MakeLowercase}[]
\titleformat{\subsection}{\large\bfseries\sffamily\uppercase}{\thesubsection}{1em}{}
\titleformat{\subsubsection}{\sffamily\bfseries}{\thesubsubsection}{1em}{}
% babel überschreibt die Namen bei \begin{document}, daher erst dort setzen
% (\ifthenelse außerhalb, weil \MakeLowercase es in der Überschrift nicht verträgt)
\ifthenelse{\boolean{english}}
  {\AtBeginDocument{%
    \renewcommand\contentsname{table of contents.}%
    \renewcommand\listfigurename{list of figures.}%
    \renewcommand\listtablename{list of tables.}%
    \renewcommand\refname{\litname.}}}
  {\AtBeginDocument{%
    \renewcommand\contentsname{inhaltsverzeichnis.}%
    \renewcommand\listfigurename{abbildungsverzeichnis.}%
    \renewcommand\listtablename{tabellenverzeichnis.}%
    \renewcommand\refname{\litname.}}}
\usepackage{exscale,relsize}
\usepackage{fancyhdr}
\usepackage[small]{caption}
\usepackage{float}
\usepackage{amsmath}
\usepackage{amssymb}
\usepackage{rotating}
\usepackage{mathrsfs}
\usepackage{booktabs}
\usepackage{listings}
\setlength{\emergencystretch}{2em}

\newcommand{\babel}[2]{\ifthenelse{\boolean{english}}{#1}{#2}}

% auxillary symbols
\newcommand{\define}{\stackrel{!}{=}}
\renewcommand{\equiv}{\,\widehat{=}\,}
\newcommand{\subsubsubsection}{\textbf}
\newcommand{\mif}{\quad\mathrm{\babel{if}{falls}}\quad}
\newcommand{\with}{\quad\mathrm{\babel{with}{mit}}\quad}
\newcommand{\for}{\quad\mathrm{\babel{for}{für}}\quad}
\newcommand{\matlab}{MATLAB}

% format specifications
\renewcommand{\emph}{\textbf}
\newcommand{\file}{\textit}
\newcommand{\cmd}{\texttt}
\newcommand{\ten}{\boldsymbol}
\newcommand{\lemma}{\textit}

% Pfade zu Grafiken
\graphicspath{{abbildungen/}{../Abgabe_Lab1/figures/}}

% mathematical operators
\newcommand{\grad}{\,\mathrm{grad}\,}
\renewcommand{\div}{\,\mathrm{div}\,}
\newcommand{\rot}{\,\mathrm{rot}\,}
\newcommand{\lap}{\Delta}
\newcommand{\laplace}[1]{\mathscr{L}\left\{#1\right\}}
\newcommand{\trans}{^T}

\setlength{\parindent}{0em}
\setlength{\parskip}{1.5ex plus0.5ex minus0.5ex}
\setlength{\captionmargin}{3em}

% Quellcode (für Ausschnitte und Anhang)
\lstset{basicstyle=\ttfamily\footnotesize,frame=single,columns=fullflexible,
        keywordstyle=\bfseries,language=C++,breaklines=true}

% Offene Stellen rot markieren -- vor der Abgabe müssen alle weg sein
\newcommand{\todo}[1]{\textcolor{red}{\textbf{[TODO:} #1\textbf{]}}}
